\section{On policy distillation with precision matched rollout}
\label{sec:C08:frontier}
The DeepSeek-V4 report's Section 5 separates specialist training from a later on-policy distillation merge, replacing its mixed-RL stage with that merge. It reports full-vocabulary logit distillation rather than only sampled token-level estimates and uses native quantized weights during rollout~\cite{rewrite_deepseek4}. These current choices connect training objectives to serving numerics; they should not be summarized as a claim that one classical preference method describes the entire pipeline.

For a fixed prefix, teacher distribution $q$, and student distribution $p_\theta$, reverse KL is $D(p_\theta\|q)=\sum_vp_\theta(v)\log[p_\theta(v)/q(v)]$. Its gradient is
\begin{equation}
 \nabla_\theta D(p_\theta\|q)=
 \mathbb E_{v\sim p_\theta}
 \left[\log\frac{p_\theta(v)}{q(v)}
 \nabla_\theta\log p_\theta(v)\right],
\end{equation}
assuming a finite vocabulary, differentiability, and $q(v)>0$ wherever $p_\theta(v)>0$. The constant score term vanishes because $\sum_v\nabla p_\theta(v)=0$. Summing over the vocabulary evaluates this conditional expectation directly; a sampled estimator adds conditional sampling variance. Prefixes generated by the student remain random, so full-vocabulary evaluation does not make the entire training gradient deterministic. Sampling tokens from the teacher instead without correction estimates a different expectation.

The conditioning distribution also matters. Let a finite rollout horizon and a fixed, parameter-independent rule for selecting a prefix position define the normalized distribution $d_\theta(h)$ of its student-generated prefix $h$ and let $K_\theta(h)=D(p_\theta(\cdot\mid h)\|q(\cdot\mid h))$. For $J(\theta)=\sum_h d_\theta(h)K_\theta(h)$, differentiating the finite sum gives
\begin{equation}
 \nabla J=\mathbb E_{h\sim d_\theta}
 \left[\nabla K_\theta(h)
 +K_\theta(h)\nabla\log d_\theta(h)\right].
\end{equation}
Holding sampled prefixes fixed while differentiating the vocabulary loss estimates the first term. It is a conditional distillation update, rather than an unbiased estimator of the displayed total derivative, unless the second term is included or vanishes under additional assumptions. This deduction identifies a reproducibility question; it does not assert an undocumented choice in the source implementation. Record whether rollout is detached, whether prefix score terms are used, and which objective the reported gradient is intended to optimize.

The memory cost can be substantial. In a hypothetical batch with eight sequences, 4096 positions, vocabulary size 128,000, and four-byte logits, one full logits tensor occupies 15.625 GiB. Holding both teacher and student tensors would occupy 31.25 GiB before activations or gradients. Chunked vocabulary reductions, fused loss kernels, and selective position processing can change the residency pattern. A study should identify which optimization is used and verify that it preserves the stated KL direction and normalization. Avoiding materialization is not the same as replacing the full objective with one sampled token.

Precision consistency creates a second falsifiable invariant. If rollout uses quantized weights and training recomputes probabilities with a different effective model, the recorded behavior distribution and optimized distribution can disagree before learning begins. Quantization-aware training can reduce that mismatch, but correctness still depends on scales, dequantization, kernel arithmetic, and sampling transforms. Compare logits and action probabilities under rollout and training paths on identical live contexts, stratified by sequence length and rare tool arguments. A small mean discrepancy does not rule out failures concentrated at an action boundary.

Distillation can also transfer the teacher's blind spots and erase useful diversity among specialists. Evaluate each domain separately before and after merging, including tool-use failures and memory conflicts, rather than accepting a higher aggregate score as evidence of uniform improvement. Keep a fixed external verifier and an untouched task set. An online student can visit prefixes absent from the original specialist data, which is one motivation for on-policy distillation, but teacher predictions there are not guaranteed correct. The appropriate evidence is verified downstream behavior and retained regressions, not the mere availability of dense teacher logits.

Preference-loss implementations also require an exact boundary between prompt and completion tokens. If the common prompt contributes to one candidate's summed log probability but not the other's, the optimized pair no longer represents the intended response comparison. Tokenization at the concatenation boundary can change the first completion token, so independently tokenizing prompt and answer is not always equivalent to tokenizing their combined serialization. Check these boundaries with actual tokenizer artifacts, including whitespace and special messages. Compare summed conditional log probabilities against a direct teacher-forced reference on a small deterministic batch before interpreting any training gain.

\paragraph{Sequence probabilities and length normalization.}
The preference derivation developed below uses complete-response probabilities, hence sums of conditional token log probabilities. Replacing those sums with per-token averages changes the objective and generally loses the original reward reparameterization. A concrete hypothetical pair shows the difference. Suppose the preferred ten-token response has student probability 0.1 and reference probability 0.01, while a dispreferred two-token response has probabilities 0.2 and 0.1. The sequence log-ratio difference is $\log 10-\log2=\log5$, favoring the preferred response despite its lower absolute student probability. Dividing each log ratio by its response length gives $\log10/10-\log2/2<0$, reversing the ordering.

These numbers are a calculated counterexample, not model measurements. They establish that length normalization is a substantive modeling choice. It can be useful for some objectives, but cannot be introduced as an innocuous numerical stabilization while retaining every theoretical claim of the original sequence objective. The same care applies to masking an end-of-sequence token or truncating long candidates. A completion that is cut off before its action or answer is not simply a shorter observation of the same preference event. The trainer should record truncation policy and the fraction of affected pairs.

An executable verification can construct a tiny vocabulary model with manually specified logits and enumerate every completion in a bounded response space. Compute the exact pair loss and gradient by direct probability products, then compare the trainer's token gathering, attention masks, reference scores, and reduction. Repeat with unequal response lengths, identical prefixes, padded examples, and an empty supervised suffix. This checks implementation against the chosen objective independently of whether a large training run happens to improve a benchmark. It is particularly valuable when the trainer supports multiple loss variants behind one public class name.

Tool trajectories introduce another reduction choice. A long episode can contain many observation tokens that are not policy actions. Including them in the policy ratio can assign optimization pressure to a distribution the agent does not control. Excluding them while accidentally retaining their padding or shifted positions can corrupt the action-token alignment. The trace should carry explicit action masks generated from the authoritative serialization, not inferred later from role-like strings in untrusted tool text. Compare the masked sum with an independently assembled action list on a small trajectory containing multiple tools and failures.

Finally, report whether learning weights episodes, decisions, or action tokens. Equal episode weighting protects short tasks from being overwhelmed by long traces, while token weighting can emphasize detailed reasoning and extensive tool interactions. Neither choice is universally correct. The relevant question is which deployment objective the weighting approximates and whether the evaluation measures that objective. A gain on a benchmark dominated by short tasks can coexist with a regression on rare long sessions if the training reduction changes their relative influence. Retaining both stratified outcomes and the exact reduction makes that tradeoff scientifically interpretable.


